\section{强化学习}
\label{sec:rl}

\gls{rl}（强化学习）是 \gls{ml} 的一个分支，旨在由 \gls{mdp} 建模的序贯决策问题中学习最优策略。强化学习的主要难点在于通过与环境交互收集数据并从中学习。\emph{样本复杂度}（智能体通过与环境交互收集的样本数）是衡量强化学习算法的重要指标。从环境中采样在时间、计算资源或资金方面可能代价高昂。虽然可以用仿真代替，但这会引入模型从仿真迁移到现实时的问题\footnote{值得注意的是，延迟在仿真中通常被忽略，在现实世界测试时可能引发问题。}。因此，常用样本复杂度来衡量强化学习算法的效率。

尽管智能体的性能仍是设计者希望优化的最终目标，但最终目标与样本复杂度之间存在一定权衡：一个学习快速的智能体可能优于保证达到最优行为但需要极长时间的智能体。这一权衡称为探索-利用问题，其中智能体必须在执行已知收益较高的动作与执行收益不确定的动作之间权衡，以获取对环境的了解。

历史上，最早应用于 \gls{mdp} 问题的方法是\emph{动态规划}算法，如\emph{策略迭代} \cite{howard1960dynamic}。但这些方法需要了解转移动力学并在 \gls{mdp} 的每个状态中维护价值函数，因此受到限制。向更现实的问题推进时，强化学习已去除了已知动力学的假设，并开始考虑更大的状态空间，最终扩展到连续空间，从而需要函数近似。本节后续部分将深入探讨其中一些解决方案，这些方案在本文中很有用，因为延迟强化学习中的许多算法都基于它们。将首先介绍针对折扣回报准则设计的算法，最后介绍针对平均奖励的算法。

\subsection{基于价值的方法}
解决强化学习问题的第一种方法是学习价值函数 $V$ 或动作价值函数 $Q$，这正是基于价值的方法所做的。开创性方法包括 SARSA \cite{rummery1994line} 和 Q学习（Q-learning）\cite{watkins1989learning}，它们基于动态规划。

\subsubsection{Q学习}
Q学习 \cite{watkins1989learning} 是一种流行算法，通过与环境交互并利用最优 Bellman 方程 \cref{eq:bellman_optimal_equation_Q} 更新估计值，递归地学习最优 Q 函数。该算法通过与环境交互采样形如 $(s_t,a_t,r_t,s_{t+1})$ 的四元组，并用于更新 Q 函数：
\begin{align}
\label{eq:update_q_learning}
    \hat Q(s_t, a_t)_i = \hat Q(s_t, a_t) + \alpha \delta_{TD,t},
\end{align}
其中 $\delta_{TD,k}$ 称为当前 Q 函数与其一步自举之间的\gls{td}（时序差分）误差。形式化地，在第 $t$ 步：
\begin{align}
\label{eq:td_error_q_learning}
    \delta_{TD,t} = r_t+\gamma \max_{a\in\actionspace} \hat Q(s_{t+1}, a) -\hat Q(s_t, a_t).
\end{align}
Q学习是一种\emph{离策略}（off-policy）算法，因为它用于从环境采样的策略与用于更新的策略不同。通常使用关于 $\hat{Q}$ 的 $\epsilon$-贪心策略进行采样，即该策略以 $1-\epsilon$ 的概率选择最大化 $\hat{Q}$ 的动作，以 $\epsilon$ 的概率在 $\actionspace$ 中均匀随机选择动作。而更新中 \gls{td} 误差项 (\Cref{eq:td_error_q_learning}) 使用的策略关于 $\hat Q$ 是贪心的，因为它始终选择最大化动作。采样的 $\epsilon$-贪心策略（或其变体）对探索很有用，它确保 \gls{mdp} 的每个状态-动作对将被无限频繁地访问，这是算法收敛性证明的要求 \cite{watkins1992q}，使智能体能评估不同动作的效果，从而可能获得更高奖励。

\subsubsection{SARSA}
SARSA \cite{rummery1994line} 算法与 Q 学习非常相似。它是一种\emph{在策略}（on-policy）算法，智能体使用当前策略从环境中采样五元组 $(s_t,a_t,r_t,s_{t+1},a_{t+1})$。使用这些样本，\gls{td} 误差表达式变为：
\begin{align*}
    \delta_{TD,t} = r_t+\gamma\hat Q(s_{t+1}, a_{t+1})-\hat Q(s_t, a_t).
\end{align*}
注意，$a_{t+1}$ 替换了 \Cref{eq:td_error_q_learning} 更新中最大化 Q 的动作。

\subsubsection{深度Q网络}
当状态空间 $\statespace$ 过大时，一种解决方案是对 Q 函数采用近似方法。一种选择是使用\emph{深度神经网络}作为近似器，如 \gls{dqn}（深度Q网络）\cite{mnih2013playing}。注意动作空间 $\actionspace$ 仍保持离散，因此 \Cref{eq:td_error_q_learning} 中的最大值可计算。放宽离散动作空间假设的工作已由其他方法提出，如深度确定性策略梯度（DDPG）\cite{lillicrap2015continuous}。

在 \gls{dqn} 中，Q 函数由参数 $\boldsymbol\theta$ 参数化，目标是找到 $Q_{\boldsymbol\theta}(s,a) \sim Q^\star(s,a)$。给定当前参数集 $\boldsymbol\theta_i$，下一参数集 $\boldsymbol\theta_{i+1}$ 通过最小化以下 \gls{td} 误差获得：
\begin{align}
\label{eq:update_dqn}
    \delta_{TD}(\theta_i,\rho) = \expectedvalue_{(s,a,r,s')\sim\rho}\left[\left(r+\gamma\max_{a\in\actionspace}Q_{\boldsymbol\theta_{i}}(s') - Q_{\boldsymbol\theta_{i+1}}(s,a)\right)^2\right].
\end{align}
其中 $\rho$ 是 $\statespace\times\actionspace\times[0,\maximumreward]\times\statespace$ 上的分布。通常在 \Cref{eq:update_dqn} 最小化的每次迭代中，使用关于 $Q_{\boldsymbol\theta_{i+1}}$ 的 $\epsilon$-贪心策略收集新样本 $(s,a,r,s')$。\cite{mnih2013playing} 观察到，在回放缓冲区中不仅保留最近采样的四元组，还保留较旧的样本来最小化 \Cref{eq:update_dqn}，可以使学习更平滑，并使问题更接近 \gls{sl}（监督学习）。这一技术称为\emph{经验回放}（experience replay）。因此，缓冲区中填充了来自许多不同策略的样本，将这些策略的混合称为 $\rho$。

\subsection{基于策略的方法}
\label{subsec:policy_based}
基于价值的方法的一个主要限制是智能体学习到关于环境的额外知识，这些知识可能不直接有用。最终，唯一关心的是学习最优行为，而学习 Q 函数只是实现这一目标的间接方式。相比之下，基于策略的方法直接从与环境的交互中学习策略。

通常，策略被视为参数化函数，从而将问题从函数集合上的搜索转化为参数集合上的搜索。这一框架称为 \gls{po}（策略优化）\cite{deisenroth2013survey}，将策略置于参数化集合 $\Pi_{\Theta} = \{\pi_{\vectorialform{\theta}} : \vectorialform{\theta} \in \Theta \subseteq \realnumbers^{n_\theta}\}$ 中考虑。

策略完全由其参数 $\vectorialform{\theta}$ 指定，因此简化其回报符号为 $\expectedreturn[][](\vectorialform{\theta})=\expectedreturn[\gamma][\pi_\vectorialform{\theta}]$。智能体的新目标变为：
\begin{align}
\label{eq:objective_policy_opt}
	\vectorialform{\theta}^\star \in \argmax_{\vectorialform{\theta} \in \Theta} \expectedreturn[][](\vectorialform{\theta}).
\end{align}

策略 $\pi_{\vectorialform{\theta}}$ 需要是随机化的，以提供足够的探索。若策略对 $\vectorialform{\theta}$ 可微，则策略优化的以下基本结论成立。

\begin{thm}[策略梯度定理]
    设 $\pi_{\vectorialform{\theta}}\in\Pi_{\Theta}$ 是随机化策略，对 $\vectorialform{\theta}$ 可微，则其回报的梯度为
    \begin{align*}
        \nabla_{\vectorialform{\theta}} \expectedreturn[][](\vectorialform{\theta}) =  \frac{1}{1-\gamma}\expectedvalue_{\substack{s\sim\discountedstateoccupancydistribution[\gamma][\pi_{\vectorialform{\theta}}]\\a\sim\pi_{\vectorialform{\theta}}(\cdot\vert s)}}\left[Q^{\pi_{\vectorialform{\theta}}}(s,a)\nabla_{\vectorialform{\theta}} \log \pi_{\vectorialform{\theta}}(a\vert s )\right]
    \end{align*}
\end{thm}
注意，许多工作（包括最先进的方法）使用梯度的近似，从状态分布中去除折扣，但保留在 Q 函数中。该量可能不再是任何函数的梯度 \cite{nota2020policy}，但提供了良好的实证结果。

策略梯度不需要访问 \gls{mdp} 模型，可以从样本中估计。以下两段介绍两种主要估计器，它们的公式假设可以访问轨迹数据集 $(\tau_i)_{1\le i\le n}$，其中每条轨迹长度为 $H$，具体为 $\tau_i=(s_0^i,a_0^i,r_0^i\dots,s_{H}^i)$。

\subsubsection{REINFORCE}
\cite{williams1992simple} 提出的 REINFORCE 估计器为：
\begin{align*}
    \widehat{\nabla_{\vectorialform{\theta}} \expectedreturn[][](\vectorialform{\theta})} = \frac{1}{n}\sum_{i=1}^{n}\sum_{t=1}^{H-1} \nabla_{\vectorialform{\theta}} \log \pi_{\vectorialform{\theta}}(a_t^i\vert s_t^i)\left(\sum_{t'=0}^{H-1}\gamma^{t'} r_{t'}^i\right).
\end{align*}
该估计器通常方差过高，但可以从最右边项中减去常数基线 $b$，而不引入偏差，同时允许控制方差 \cite{williams1992simple}。

\subsubsection{策略梯度定理}
降低上述估计器方差的另一种方法是注意到未来动作不影响过去奖励。实际上，对于任意 $i>j$，记 $\discountedstateoccupancydistribution[\gamma,j:i][\pi]$ 为策略 $\pi$ 下轨迹片段 $(s_j,a_j,\dots,s_i,a_{i})$ 的分布，则有：
\begin{align*}
    &\expectedvalue_{(s_j,a_j,\dots,s_i,a_{i})\sim\discountedstateoccupancydistribution[\gamma,j:i][\pi_{\vectorialform{\theta}}]}\left[\nabla_{\vectorialform{\theta}} \log \pi_{\vectorialform{\theta}}(a_i\vert s_i)r(s_{j},a_{j})\right]
    \\&\quad= \expectedvalue_{\substack{(s_j,a_j,\dots,s_{i-1},a_{i-1})\sim\discountedstateoccupancydistribution[\gamma,j:i-1][\pi_{\vectorialform{\theta}}]\\s_i\sim p(\cdot\vert s_{i-1},a_{i-1})}}\left[ r(s_{j},a_{j})\int_{\actionspace}\pi_{\vectorialform{\theta}}(a_i) \nabla_{\vectorialform{\theta}} \log \pi_{\vectorialform{\theta}}(a_i\vert s_i)\;\de a_i\;\right]
    \\&\quad= \expectedvalue_{\substack{(s_j,a_j,\dots,s_{i-1},a_{i-1})\sim\discountedstateoccupancydistribution[\gamma,j:i-1][\pi_{\vectorialform{\theta}}]\\s_i\sim p(\cdot\vert s_{i-1},a_{i-1})}}\left[r(s_{j},a_{j})\int_{\actionspace}\nabla_{\vectorialform{\theta}}\pi_{\vectorialform{\theta}}(a_i) \;\de a_i\; \right]
    \\&\quad= \expectedvalue_{\substack{(s_j,a_j,\dots,s_{i-1},a_{i-1})\sim\discountedstateoccupancydistribution[\gamma,j:i-1][\pi_{\vectorialform{\theta}}]\\s_i\sim p(\cdot\vert s_{i-1},a_{i-1})}}\left[r(s_{j},a_{j})\nabla_{\vectorialform{\theta}}1\right]
    \\&\quad= 0
\end{align*}
由此得到策略梯度定理（PGT）\cite{sutton1999policy}：
\begin{align}
\label{eq:pgt}
    \widehat{\nabla_{\vectorialform{\theta}} \expectedreturn[][](\vectorialform{\theta})} = \frac{1}{n}\sum_{i=1}^{n}\sum_{t=1}^{H-1} \nabla_{\vectorialform{\theta}} \log \pi_{\vectorialform{\theta}}(a_t^i\vert s_t^i)\left(\sum_{t'=t}^{H-1}\gamma^{t'} r_{t'}^i-b(s_t)\right).
\end{align}
注意最右边项的和现在从 $t$ 开始。与 REINFORCE 一样，此处也引入了基线 $b$ 以进一步降低方差。在 PGT 中，基线可依赖于当前状态 $s_t$ 而不引入偏差 \cite{peters2008reinforcement}。\cite{peters2008reinforcement} 讨论了基线选取的策略。

\subsubsection{信任域策略优化}
\gls{trpo}（信任域策略优化）\cite{schulman2015trust} 是一种基于策略的算法，其思想是提供安全的改进步骤，即保证性能提升的策略更新。其理论基础依赖于著名的性能差异引理，在此回顾。

\begin{lemma}[性能差异引理，\cite{kakade2002approximately} 的引理 6.1]
    设 $\pi,\pi'\in\Pi^{SR}$，其期望折扣回报分别为 $\expectedreturn[\gamma][\pi]$ 和 $\expectedreturn[\gamma][\pi']$，则
    \begin{align*}
        \expectedreturn[][\pi'] - \expectedreturn[][\pi] = \frac{1}{1-\gamma}\expectedvalue_{(s,a)\sim\discountedstateoccupancydistribution[\gamma][\pi']}\left[\advantagefunction^{\pi}(s,a)\right].
    \end{align*}
\end{lemma}

基于该引理，\cite{schulman2015trust} 构建了一个替代目标函数，其最大化保证了 $\expectedreturn[\gamma][\pi']$ 的提升。

令 $\textstyle C=\frac{4\varepsilon\gamma}{(1-\gamma)^2}$ 和 $\varepsilon=\max_{s,a}\advantagefunction^\pi(s,a)$，\cite{schulman2015trust} 证明：
\begin{align*}
    \expectedreturn[\gamma][\pi'] \ge \underbrace{\expectedreturn[\gamma][\pi] + \frac{1}{1-\gamma}\expectedvalue_{(s,a)\sim\discountedstateoccupancydistribution[\gamma]}\left[\advantagefunction^{\pi'}(s,a)\right]}_{\coloneq L^\pi(\pi')} - C \max_s \kullbackleiblerdivergence^{\max} (\pi(\cdot\vert s),\pi'(\cdot\vert s)).
\end{align*}

为简化优化，\cite{schulman2015trust} 在最终目标中用平均 KL 散度替代最大 KL 散度，前者在实践中更易计算。此外，$C$ 中的惩罚项被替换为约束，以允许更大但稳健的更新步长。最大化问题变为：
\begin{align*}
\begin{array}{cl}
     \underset{\pi'}{\text{maximise}} &  L^\pi(\pi')\\
     \text{subject to} & \expectedvalue_{s\sim \discountedstateoccupancydistribution[\gamma][\pi]}\left[ \kullbackleiblerdivergence (\pi(\cdot\vert s),\pi'(\cdot\vert s))\right]\le \delta.
\end{array}
\end{align*}
\gls{trpo} 已被证明是一种经验高效的算法。

\subsubsection{近端策略优化}
尽管 \gls{trpo} 取得了实证成功，但其计算成本较高，因为需要计算 KL 散度的 Hessian 矩阵及其逆。相比之下，\gls{ppo}（近端策略优化）\cite{schulman2017proximal} 提出使用裁剪操作将问题还原为一阶优化方案，以保持策略在安全区域内。实证结果表明，\gls{ppo} 与 \gls{trpo} 至少同样高效，尤其在计算更简单方面表现更优。

\subsubsection{基于参数的优化}
\label{subsec:param_based_optim}
从参数化 \gls{po} 出发，可以增加一层抽象，考虑优化\emph{超策略}而非策略本身。基本上，超策略定义了智能体将遵循的策略的选择规则。在此设定中，策略的参数 $\vectorialform{\theta}$ 由超策略 $\nu_{\vectorialform{\rho}}$ 采样，超策略本身由参数向量 $\vectorialform{\rho}$ 参数化，其中 $\vectorialform{\rho}$ 属于 $\mathcal{V}_{\mathcal{P}} = \{\nu_{\vectorialform{\rho}} : \vectorialform{\rho} \in \mathcal{P} \subseteq \realnumbers^{n_\rho}\}$ \cite{sehnke2008parameter}。将随机性转移到超策略后，策略不再需要是随机化的。这对于降低估计器方差是一个巨大优势。实际上，只需采样单一 $\vectorialform{\theta}$ 即可收集一条轨迹，轨迹的随机性仅来自环境。此框架中的目标变为：
\begin{align}
\label{eq:objective_hyperpolicy_opt}
	\vectorialform{\rho}^\star \in \argmax_{\vectorialform{\rho} \in \mathcal{P}} \expectedreturn[][](\vectorialform{\rho})=\argmax_{\vectorialform{\rho}\in \mathcal{P}} \expectedvalue_{\boldsymbol\theta\sim \nu_{\vectorialform{\rho}}}\expectedreturn[][](\vectorialform{\theta}).
\end{align}
将此设定称为基于参数的 \gls{po}，将原始设定称为基于动作的 \gls{po}。

\subsection{Actor-Critic}
在 \Cref{eq:pgt} 的公式中，项 $\textstyle\sum\nolimits_{t'=t}^{H-1}\gamma^{t'} r_{t'}$ 实质上是 Q 函数 $Q(s_t,a_t)$ 的蒙特卡洛估计。\emph{Actor-Critic} 方法提出将该估计替换为 \gls{td} 估计，如 TD$(0)$：
\begin{align*}
    G_{t}^0 = r_{t} + \gamma V_{\boldsymbol\phi}(s_{t+1}).
\end{align*}
这里 $V_{\boldsymbol\phi}$ 是使用参数 $\boldsymbol\phi$ 对状态价值函数的当前近似。在 Advantage Actor-Critic (A2C) 算法 \cite{mnih2016asynchronous} 中，选择 $V_{\boldsymbol\phi}(s_{t})$ 作为 \Cref{eq:pgt} 中的基线是一个有趣的选择。它将 \Cref{eq:pgt} 中最右边的括号项变为：
\begin{align}
\label{eq:a2c}
    \widehat{\nabla_{\vectorialform{\theta}} \expectedreturn[][](\vectorialform{\theta})} = \frac{1}{n}\sum_{i=1}^{n}\sum_{t=1}^{H-1} \nabla_{\vectorialform{\theta}} \log \pi_{\vectorialform{\theta}}(a_t^i\vert s_t^i)
    \left(\vphantom{\sum_{i=1}^{n}\sum_{t=1}^{H-1}} r_{t} + \gamma V_{\boldsymbol\phi}(s_{t+1}) - V_{\boldsymbol\phi}(s_{t})\right).
\end{align}
注意新项对应于 $r_{t} + \gamma V_{\boldsymbol\phi}(s_{t+1}) - V_{\boldsymbol\phi}(s_{t})=\advantagefunction_{\boldsymbol\phi}(s_t,a_t)$，这是优势函数的估计（参见\Cref{eq:advantage_function}）。Actor-Critic 的名称源于该方法同时学习策略 $\pi_{\vectorialform{\theta}}$（\emph{actor}）和状态价值函数 $V_{\boldsymbol\phi}$（\emph{critic}）。该算法交替使用 \Cref{eq:a2c} 的梯度估计更新策略，并使用以下目标最小化期望平方 \gls{td} 误差来更新价值函数（参见 \Cref{eq:update_dqn} 作比较）：
\begin{align}
\label{eq:update_a2c}
    \delta_{TD}(\boldsymbol\phi_i,\rho) = \expectedvalue_{(s,a,r,s')\sim\rho}\left[\left(r+\gamma V_{\boldsymbol\phi_i}(s') - V_{\boldsymbol\phi_{i}}(s)\right)^2\right].
\end{align}
同样，$\rho$ 是某个转移数据集 $\mathcal D$ 上的分布。Actor-Critic 相对于 PGT 和 REINFORCE 的一个主要优势是它支持离策略学习。

\subsubsection{Soft Actor-Critic}
本文中将使用的 Actor-Critic 算法示例是 \gls{sac}（Soft Actor-Critic）\cite{haarnoja2018soft}。其主要特点是在上述标准 A2C 的更新中增加了熵正则化。对于随机变量 $X$，其概率为 $p(x)=\probability(X=x)$，熵定义为：
\begin{align*}
    \entropy(X) = \expectedvalue_{x\sim p}[-\log p(x)].
\end{align*}
熵正则化 \gls{rl} 在原始奖励函数中增加了策略熵的奖励。这意味着回报现在变为：
\begin{align}
\label{eq:discounted_reward_entropy}
    \expectedreturn[\gamma] = (1-\gamma)\expectedvalue\left[\sum_{t=1}^H \gamma^t  \probability\left(S_t=s; \pi,\mu\right)\pi(a\vert s) \left(r(s,a)v + \alpha \entropy(\pi(\cdot\vert s))\right)\right],
\end{align}
其中 $\alpha>0$ 是控制正则化强度的参数。

根据这一新的回报，可提出新的状态价值函数和状态-动作价值函数，\gls{td} 误差也与 \Cref{eq:update_a2c} 不同。因此，在 \gls{sac} 中，状态-动作价值函数通过最小化以下目标学习：
\begin{align}
\label{eq:update_sac}
    \delta_{TD}(\boldsymbol\phi_i,\rho) &= \quad\expectedvalue_{(s,a,r,s')\sim\rho}\left[\left(y(r,s') - Q_{\boldsymbol\phi_{i}}(s,a)\right)^2\right],
\end{align}
其中
\begin{align*}
    y(r,s') = \expectedvalue_{a'\sim\pi_{\boldsymbol\theta_{i}}}\left[r+\gamma \min_{j\in\{1,2\}}Q_{\boldsymbol\phi_{\text{target,j}}}(s',a')-\alpha \log\pi_{\boldsymbol\theta_{i}}(a'\vert s') \right].
\end{align*}
目标 Q 函数 $(Q_{\boldsymbol\phi_{\text{target,j}}})_{j\in\{1,2\}}$ 是最近几次迭代中 Q 函数参数的 Polyak 平均值，用于替代当前 Q 函数 $Q_{\boldsymbol\phi_i}$，以避免实证中观察到的不稳定性。$\textstyle\min_{j\in\{1,2\}}$ 是另一个避免过高估计的经验性技巧。完整算法参见 \cite{haarnoja2018soft}。

\subsection{模仿学习}
\label{subsec:imitation_learning}
\emph{模仿学习}这一范式源于以下观察：从示范中学习技能通常比从零开始学习更容易。实际上，在固定的观测集上学习模仿专家策略属于\emph{\glsfirst{sl}}（监督学习）框架，通常比 \gls{rl} 更简单。因此，模仿学习旨在缩小 \gls{rl} 与 \gls{sl} 之间的差距。为更形式化地定义该框架，设 $\pi_E$ 为专家策略；大多数模仿学习方法的目标是找到学习者的策略 $\pi_I$，使得 $\textstyle\expectedvalue_{s\sim d_\mu^{\pi_E}}[l(s,\pi_I)]$ 最小化 \cite{ross2011reduction}。函数 $l(s,\pi)$ 是一个损失函数，旨在使 $\pi_I$ 更接近 $\pi_E$。值得注意的是，该目标是在 $\pi_E$ 诱导的状态分布下定义的。这可能带来问题，因为一旦学习者犯错，它可能进入一个对专家行为了解不足的状态，从而进一步犯错。因此，误差可能按有效水平线的平方传播 \cite[Theorem~2.1]{ross2010efficient}\cite[Theorem~1]{xu2020error}。

其他模仿学习方法旨在解决这一问题，但通常不实用，如 \cite{ross2011reduction} 所述。例如，\cite{ross2010efficient} 提出使用非平稳策略，这需要为每个时间步训练一个策略，对长水平线而言计算成本高昂。\cite{ross2010efficient} 提出的另一种方法 SMILe 则考虑策略的混合，每次迭代增加一个策略。如 \cite{ross2011reduction} 所述，这在实践中存在问题，因为混合中包含不同质量的策略，较弱策略可能犯错，而较强策略必须弥补这些错误，这显然会造成不稳定。

一个相当成功的解决方案称为数据集聚合（\textsc{DAgger}）\cite{ross2011reduction}。该模仿算法在学习者的状态分布下计算损失 \cite{osa2018algorithmic}，因此能够考虑学习者策略可能犯错所导致的分布偏移。后文将看到，在延迟情况下，延迟会加剧延迟策略与无延迟策略之间状态分布的差异，因此这是模仿算法所需的一个理想性质。实际上，\textsc{DAgger} 的思想很简单：根据与学习者相似的策略收集新样本，然后仅在这些样本上查询专家，了解专家本应采取的动作。这提供了与学习者状态分布相匹配的样本，是一个巨大优势。采样策略为 $\pi_i = \beta_i \pi_E + (1-\beta_i) \hat{\pi}_i$，其中 $\beta_i$ 是专家策略与学习者当前策略 $\hat{\pi}_i$ 混合的权重。通过将采样状态 $s$ 与专家选择的动作 $\pi_E(s)$ 相结合，构建用于 \gls{sl} 步骤的数据集 $\mathcal{D}$。然后在 $\mathcal{D}$ 上训练新的模仿策略 $\hat{\pi}_{i+1}$。序列 $(\beta_i)_{i\in[\![1,N]\!]}$ 满足 $\beta_1=1$（初始仅从 $\pi_E$ 采样）和 $\beta_N=0$（最终仅从模仿策略采样）。

\subsection{理论强化学习}
\label{subsec:theoretical_rl}

在最后一小节中，介绍针对平均奖励准则设计的方法。这些方法通常更具理论性，关注平衡 \gls{mdp} 的探索与已获取知识的利用。智能体必须学习在这两种策略之间权衡，以免相对于专家损失太多。损失可以理解为因缺乏探索而错失机会，或因缺乏利用而损失奖励。该设定属于\emph{在线学习}范畴。为更好理解该设定，介绍在线学习的一个子领域：\emph{个体序列预测} \cite{cesa2006prediction}。该领域考虑一个智能体需要在 $K$ 个动作（称为臂）之间反复选择。每个臂为智能体提供奖励，奖励可以随机方式或对抗方式产生。环境没有状态，智能体可随时在 $K$ 个臂之间选择。该框架可表示为具有单一状态的 \gls{mdp}，即 $\cardinal{\statespace}=1$。当智能体只能观察到所选臂的奖励，而无法观察到其他 $K-1$ 个臂的奖励时，该设定称为\emph{多臂赌博机} \cite{lattimore2020bandit}。当奖励是随机时，称为\emph{随机赌博机}；当奖励是对抗性时，称为\emph{对抗赌博机}。

在该领域中，一个重要的概念是\emph{遗憾}（regret），即学习智能体收集的奖励与事后最优策略收集的奖励之间的期望差异。目标是获得次线性（sub-linear）的遗憾，表明智能体越来越接近专家策略。解决该问题的常用方法是在不确定性面前保持\emph{乐观}。在估计量（臂的奖励）上计算上置信界，智能体选择具有最佳上界的臂。后文将在 \Cref{subsec:reward_delay_related} 中看到延迟在赌博机文献中是如何处理的。

上置信界的思想也已应用于 \gls{rl}。先驱者是 UCRL 算法 \cite{auer2006logarithmic}，它对奖励函数和转移函数维护置信界。在 \gls{rl} 中，对于某个 \gls{mdp} $\markovdecisionprocess$，遗憾定义为：
\begin{align*}
    \mathcal{R}(T,s) \coloneq T\cdot\max_{\pi\in\Pi^{\text{MR}}} \averagereturnfunction(s) - \sum_{t=1}^T \expectedrewardfunction(s_t,a_t)
\end{align*}
其中 $\averagereturnfunction(s)$ 是具有初始状态分布 $\delta_s$ 的最优马尔可夫策略的平均奖励性能，$\textstyle \sum_{t=1}^T \expectedrewardfunction(s_t,a_t)$ 是学习算法在前 $T$ 步中收集的奖励。UCRL 利用关于 \gls{mdp} 混合时间的信息并假设过程的遍历性，以实现次线性遗憾。

后文将看到，后续方法在更小的遗憾和更弱的假设方面改进了 UCRL 的结果。从这些方法中，重点关注 UCRL2 \cite{auer2008near}，它适用于\emph{连通} \glspl{mdp} \cite[Section~8.3.1]{puterman1994markov}。下面回顾连通 \glspl{mdp} 的定义。

\begin{defi}[连通 \gls{mdp}]
\label{def:communicating_mdp}
    \gls{mdp} $\markovdecisionprocess$ 称为连通的，若对于任意状态 $s$ 和 $s'$，存在一个确定性平稳策略，使从 $s$ 出发能以非零概率到达 $s'$。
\end{defi}

注意在\emph{连通} \glspl{mdp} 中，平均奖励 $\averagereturnfunction(s)$ 不再依赖于 $s$。UCRL2 的结果依赖于 \gls{mdp} 的\emph{直径}（diameter）概念，该概念与连通性密切相关。事实上，有限直径与连通 \glspl{mdp} 是等价的 \cite{auer2008near}。其定义如下。

\begin{defi}[\gls{mdp} 的直径 \cite{auer2008near}]
\label{def:diameter}
    设 $\markovdecisionprocess$ 为一个 \gls{mdp}，$\policy$ 为 $\markovdecisionprocess$ 上的一个平稳策略。令 $T(s'\vert s;\policy)$ 为测量策略 $\pi$ 从 $s$ 出发首次到达 $s'$ 所需时间的随机变量。则 $\markovdecisionprocess$ 的直径为：
    \begin{align*}
        D(\markovdecisionprocess) \coloneq \max_{s\neq s'\in\statespace}\min_{\pi\in\Pi^{\text{S}}} \expectedvalue\left[T(s'\vert s;\policy)\right].
    \end{align*}
\end{defi}

利用直径的知识，\cite{auer2008near} 提出了基于乐观主义的算法分析，证明了次线性遗憾。
